\section{Introduction}

As a reminder, we use the following shortcuts:

\begin{Luacode}
local ld = luadraw -- alias sur l'espace de nommage
local cpx = ld.cpx -- raccourci pour la classe cpx
local pt3d = ld.pt3d -- raccourci pour la classe pt3d
\end{Luacode}

\subsection{Prerequisites}

\begin{itemize}
    \item This chapter presents the use of the \luadrawenv package with the \emph{3d} global option:
\verb|\usepackage[3d]{luadraw}|.
    \item The package loads the \emph{luadraw\_graph2d.lua} module, which defines the \emph{ld.graph} class and provides the \luadrawenv environment for creating graphs in Lua. Everything said in the previous chapter (Drawing 2D) therefore applies, and is assumed to be known here.
    \item The \emph{3d} global option also allows the loading of the \emph{luadraw\_graph3d.lua} module. This also defines the \emph{ld.graph3d} class (which relies on the \emph{ld.graph} class) for 3D drawings.
\end{itemize}

\subsection{Some reminders}

\begin{itemize}
    \item Another global package option: \opt{noexec}. When this global option is mentioned, the default value of the \opt{exec} option for the \luadrawenv environment will be \false (and no longer \true).

    \item When a graph is finished, it is exported in TikZ format, so this package also loads the TikZ package and the libraries:

    \begin{itemize}
        \item\emph{patterns}
        \item\emph{plotmarks}
        \item\emph{arrows.meta}
        \item\emph{decorations.markings}
    \end{itemize}
    
    \item Graphs are created in a \luadrawenv environment, which calls \emph{luacode}, so the \textbf{Lua language} must be used in this environment.

    \item Saving the \emph{*.tkz} file: the chart is exported in TikZ format to a file (with the \emph{tkz} extension). By default, it is saved in the \emph{\_luadraw} folder, which is a subfolder of the current folder (containing the master document), but it is possible to specify a path to another subfolder, with the global option \opt{cachedir=\ldots}.

    \item The environment options are:
    
    \begin{itemize}
        \item \opt{name=\ldots}: allows you to name the resulting TikZ file. It is given a name without an extension (the extension will be automatically added; it is \emph{.tkz}). If this option is omitted, then a default name is used, which is the name of the master file followed by a number.

        \item \opt{exec=true/false}: allows you to execute or not the Lua code included in the environment. By default, this option is true, \textbf{EXCEPT} if the global option \opt{noexec} was mentioned in the preamble with the package declaration. When a complex graph that requires a lot of calculations is developed, it may be useful to add the option \opt{exec=false}; this will avoid recalculating the same graph for future compilations.

        \item \opt{auto=true/false}: allows you to automatically include or not the TikZ file in place of the \luadrawenv environment when the \opt{exec} option is \false. By default, the \opt{auto} option is \true.
    \end{itemize}
\end{itemize}


\subsection{Creating a 3D Graph}

\begin{TeXcode}
\begin{luadraw}{ name=<filename>, exec=true/false, auto=true/false }
local ld = luadraw
-- create a new graph and give it a local name
local g = ld.graph3d:new{ window3d={x1,x2,y1,y2,z1,z2 [,xscale,yscale,zscale]}, adjust2d=true/false, viewdir={30,60}, window={x1,x2,y1,y2 [,xscale,yscale]}, margin={top,right,bottom,left}, size={width,height,ratio}, bg="color", border=true/false }
-- build graph g
graph instructions in Lua language ...
-- display graph g and save it in the file <filename>.tkz
g:Show()
-- or Save only in the <filename>.tkz file
g:Save()
\end{luadraw}
\end{TeXcode}

\noindent\textbf{Important}: throughout the rest of this chapter, points in space are called \emph{3D points}, they are triplets of $\mathbf R^3$. In the \luadrawenv environment, the 3D point $(x,y,z)$ will be denoted \cmd{pt3d.M(x, y, z)} (\cmd{pt3d.M} is a function that creates and returns a 3D point).

Creation is done in a \luadrawenv environment. This creation is done on the first line inside the environment by naming the graph:

\begin{Luacode}
local ld = luadraw
local g = ld.graph3d:new{ window3d={x1,x2,y1,y2,z1,z2 [,xscale,yscale,zscale]}, adjust2d=true/false, viewdir={30,60}, window={x1,x2,y1,y2 [,xscale,yscale]}, margin={left,right,top,bottom}, size={width,height,ratio}, bg="color", border=true/false }
\end{Luacode}

The \emph{graph3d} class is defined in the \luadrawenv package using the global option \opt{3d}. This class is instantiated by invoking its constructor and giving it a name (here it's \emph{g}). This is done locally so that the graph \emph{g} thus created will no longer exist once it leaves the environment (otherwise \emph{g} would remain in memory until the end of the document).

\begin{itemize}
    \item The option \opt{window3d=\{x1,x2,y1,y2,z1,z2 \fac{,xscale,yscale,zscale}\}} defines the $\mathbf R^3$ parallelepiped corresponding to the graph: it is $[x_1;x_2]\times[y_1;y_2]\times[z_1;z_2]$, as well as the scale on the three axes: \emph{xscale}, \emph{yscale}, and \emph{zscale}, these are optional and default to $1$. The default parallelepiped is $[-5;5]\times[-5;5]\times[-5;5]$.

    \textbf{Caution}: The three scales determine the initial 3D matrix of the graph. When one of them is not $1$, this matrix is ​​not the identity matrix. If you need to change the graph matrix later, you must use the method \cmd{g:Composematrix3d()} and not \cmd{g:Setmatrix3d()}, and remember to save the initial matrix beforehand with the method \cmd{g:Savematrix()}, you can then restore it with the method \cmd{g:Restorematrix()}.

    \item The option \opt{adjust2d} indicates whether the 2D window that will contain the projection of the 3D drawing should be determined automatically (\false by default). This 2D window corresponds to the \opt{window} option.

    \item The option \opt{viewdir} is a table that defines the projection method and the two viewing angles (in degrees). By default  \opt{viewdir=\{"ortho",30,60\}} (orthographic projection). The following figure shows what these two angles correspond to.

\begin{center}
\begin{luadraw}{name=viewdir}
local ld = luadraw
local pt3d = ld.pt3d
local Origin, vecI, vecJ, vecK = pt3d.Origin, pt3d.vecI, pt3d.vecJ, pt3d.vecK 
local g = ld.graph3d:new{ size={8,8} }
local i, M = ld.cpx.I, ld.pt3d.M
local O, A = Origin, M(4,4,4)
local B, C, D, E = ld.pxy(A), ld.px(A), ld.py(A), ld.pz(A) --projected from A onto the xOy plane and onto the axes
g:Dpolyline3d( {{O,A},{-5*vecI,5*vecI},{-5*vecJ,5*vecJ},{-5*vecK,5*vecK}}, "->") -- axes
g:Dpolyline3d( {{E,A,B,O}, {C,B,D}}, "dashed")
g:Dpath3d( {C,O,B,2.5,1,"ca",O,"l","cl"}, "draw=none,fill=cyan,fill opacity=0.8") --angular sector
g:Darc3d(C,O,B,2.5,1,"->") -- arc de cercle pour theta
g:Dpath3d( {E,O,A,2.5,1,"ca",O,"l","cl"}, "draw=none,fill=cyan,fill opacity=0.8") --angular sector
g:Darc3d(E,O,A,2.5,1,"->") -- arc of a circle for phi
g:Dballdots3d(O) -- the point of origin in the form of a small sphere
g:Labelsize("footnotesize")
g:Dlabel3d(
  "$x$", 5.25*vecI,{}, "$y$", 5.25*vecJ,{}, "$z$", 5.25*vecK,{},
  "towards the observer", A, {pos="E"},
  "$O$", O, {pos="NW"},
  "$\\theta$", (B+C)/2, {pos="N", dist=0.15},
  "$\\varphi$", (A+E)/2, {pos="S",dist=0.25})
g:Dlabel("viewdir=\\{'ortho',$\\theta,\\varphi$\\} (in degrees)",-5*i,{pos="N"}) -- label 2D
g:Show()
\end{luadraw}
\captionof{figure}{Viewing Angles}\label{viewdir}
\end{center}

    \item The other options are those of the \emph{ld.graph} class, described in Chapter 1.
\end{itemize}

\paragraph{Graph construction.}

\begin{itemize}
    \item The instantiated object (\emph{g} in the example) has all the methods of the \emph{ld.graph} class, plus methods specific to 3D.
    \item The \emph{ld.graph3d} class also provides a number of mathematical functions specific to 3D.
\end{itemize}

\subsection{Affine Projection Modes}

\begin{itemize}
    \item Orthogonal Projections:
    \begin{itemize}
        \item \opt{viewdir=\{"ortho",theta,phi\}}, or \opt{viewdir=\{theta,phi\}}: orthographic projection (orthogonal projection onto the screen), this is the default projection with theta=30 and phi=60 (degrees),
        \item \opt{viewdir="xOy"}: orthogonal projection onto the $xy$ plane,
        \item \opt{viewdir="xOz"}: orthogonal projection onto the $xz$ plane,
        \item \opt{viewdir="yOz"}: orthogonal projection onto the $yz$ plane.
    \end{itemize}
    \item Non-orthogonal Projections:
    \begin{itemize}
        \item \opt{viewdir=\{"yz",k,alpha\}}: cavalier perspective on the $yz$ plane,
        \item \opt{viewdir=\{"xz",k,alpha\}}: cavalier perspective on the $xz$ plane,
        \item \opt{viewdir=\{"xy",k,alpha\}}: cavalier perspective on the $xy$ plane,
        \item \opt{viewdir="iso"}: isometric perspective.
    \end{itemize}
    The three cavalier perspectives are defined using two parameters: a positive number \emph{k} and an angle in degrees \emph{alpha}, which are highlighted in the following figure.
\end{itemize}

\begin{center}
\begin{luadraw}{name=perspectives}
local ld, cpx, pt3d = luadraw, luadraw.cpx, luadraw.pt3d
local Origin, vecI, vecJ, vecK = pt3d.Origin, pt3d.vecI, pt3d.vecJ, pt3d.vecK 
local Z, Zp, deg = cpx.Z, cpx.Zp, ld.deg

local k, alpha, r = 0.65, 60, 3
local g = ld.graph3d:new{
    window3d = {-4.5,4.5,-4.5,4.5,-4.5,4.5},
    window ={-5,5,-5,5},
    viewdir = {"yz",k,alpha},
    size={10,10},
}
g:Labelsize("footnotesize")
local draw = function()
    g:Dscene3d( g:addAxes(Origin, {arrows=1}) )
    g:Darc(1,0,Zp(1,alpha*deg),(r+0.5)*k,1,"->"); g:Ddots(r*k)
    g:Dlabel("$\\alpha$",Zp((r+1)*k,30*deg),{pos="NE"}, "$k$", r*k,{pos="SE"})
    g:Dcircle(0,r*k)
end
g:Dline(0,1); g:Dline(0,cpx.I)
-- top left
g:Saveattr(); g:Viewport(-5,0,0,5); g:Coordsystem(-4.5,5.5,-5,5)
draw()
g:Ddots3d(r*vecI); g:Dlabel3d("$x=1$",r*vecI,{pos="W",dist=0.1}); 
g:Dlabel("perspective on yz plane",Z(0.5,-5),{pos="N",node_options="fill=white"})
g:Restoreattr()
-- top right
g:Saveattr(); g:Viewport(0,5,0,5); g:Coordsystem(-4.75,5.5,-5,5)
g:Setviewdir({"xz",k,alpha})
draw()
g:Ddots3d(r*vecJ); g:Dlabel3d("$y=1$",r*vecJ,{pos="NW"}); 
g:Dlabel("perspective on xz plane",Z(0.5,-5),{pos="N",node_options="fill=white"})
g:Restoreattr()
-- bottom left
g:Saveattr(); g:Viewport(-5,0,-5,0); g:Coordsystem(-4.5,5.5,-5,5.5)
g:Setviewdir({"xy",k,alpha})
draw()
g:Ddots3d(r*vecK); g:Dlabel3d("$z=1$",r*vecK,{pos="W"}); 
g:Dlabel("perspective on xy plane",Z(0.5,-5),{pos="N",node_options="fill=white"})
g:Restoreattr()
-- bottom right
g:Saveattr(); g:Viewport(0,5,-5,0); g:Coordsystem(-4.5,5.5,-5,5.5)
g:Setviewdir("iso")
g:Dscene3d( g:addAxes(Origin, {arrows=1}) ); g:Dcircle(0, cpx.abs(g:Proj3d(r*vecI)))
g:Ddots3d({r*vecI,r*vecJ,r*vecK}); 
g:Dlabel3d("$z=1$",r*vecK,{pos="NW"}, "$x=1$",r*vecI,{pos="NW"},"$y=1$",r*vecJ,{pos="NE"});
g:Dlabel("isometric perspective",Z(0.5,-5),{pos="N",node_options="fill=white"})
g:Restoreattr()
g:Show()
\end{luadraw}
\captionof{figure}{Affine Projection Modes}
\end{center}

All values ​​that can be passed to the \opt{viewdir} option during creation can also be used in the \cmd{g:Setviewdir()} method during graphing.


\subsection{Central Projection}

Since version 2.4, \luadrawenv also offers central projection. Unlike previous modes, \textbf{this projection is not affine}, and furthermore, it is not defined for all points in space, which can lead to errors, thus requiring thought and adjustments. This projection is defined by:

\begin{itemize}
    \item A camera, which is a point in space stored in a global variable called \varglob{ld.camera} and which should not be modified directly.
    \item A target, which is a point in space stored in a global variable called \varglob{ld.target} and which should not be modified directly.
\end{itemize}


The plane passing through the \varglob{ld.target} and orthogonal to the \varglob{ld.target} - \varglob{ld.camera} axis is the projection plane; it represents the screen. As with the previous modes, the central projection is accessed via the \opt{viewdir} option, or the \cmd{g:Setviewdir} method:
\cmdln{ viewdir = \{"central" \fac{, camera, target}\}}
or
\cmdln{viewdir = \{"central" \fac{, theta, phi, d, target}\}}

In the first case, the values ​​of \argu{camera} and \argu{target} are given (3D points; by default, \argu{target} is the origin). In the second case, the three arguments \argu{theta}, \argu{phi}, and \argu{d} are used to calculate the position of the camera according to the following diagram:

\begin{center}
\begin{luadraw}{name=central_perspective}
local ld, cpx, pt3d = luadraw, luadraw.cpx, luadraw.pt3d
local Origin, vecI, vecJ, vecK, M = pt3d.Origin, pt3d.vecI, pt3d.vecJ, pt3d.vecK, pt3d.M

local g = ld.graph3d:new{window3d={-5,3,-5,3,-5,5}, window={-5,10,-7,7}, size={12,12}, margin={0,0,0,0}, bbox=false,
    viewdir={"central",-60,65,35}}
    
g:Writeln("\\tikzset{->-/.style= {decoration={markings, mark=at position #1 with {\\arrow{stealth}}}, postaction={decorate}}}")
g:Labelsize("footnotesize")
local dcamera = function(pos,dir)
    local a, b, c, d, e = 0.25, 0.25, 0.1, 0.2, 0.1
    local u = cpx.normalize(dir)
    local v = cpx.I*u
    local chem, dep = {}, pos+e*u-b*v
    chem = {dep,dep+2*a*u,dep+2*a*u+2*b*v,dep+2*b*v,0.15,"cla", dep+(b-c)*v,"m",pos-d*v,pos+d*v,dep+(b+c)*v,"l"}
    g:Dpath(chem)
end
local d = 10
local N = pt3d.normalize(M(2,1,1.5))
local v = pt3d.normalize(pt3d.prod(N,vecJ))
local u = pt3d.prod(v,N)
local Cam = d*N
local A = ld.pxy(Cam)
local B, C = M(0,2,2.5), M(3,-1,4)
local E = (B+C)/2+vecK
local F = pt3d.isobar3d({B,C,E})+1.51*vecJ
local T = ld.tetra(B,E-B,C-B,F-B)
local B1, C1, E1, F1 = ld.proj3dO(B, {Origin,N}, Cam-B), ld.proj3dO(C, {Origin,N}, Cam-C), ld.proj3dO(E, {Origin,N}, Cam-E), ld.proj3dO(F, {Origin,N}, Cam-F)
local x, y = 12,8
local D = Origin -x*u/2-y*v/2.5
local plan = {D,D+x*u,D+x*u+y*v,D+y*v}  
g:Dscene3d(
    g:addFacet(plan, {color="white", contrast=0.125, edge=true}),
    g:addPoly(T, {color="white", contrast=0.25, edge=true, hidden=true, hiddenstyle="dashed"}),
    g:addPolyline( {{(Cam+B)/2,B1},{C,C1},{E,E1},{F,F1}}, {width=2,color="gray", hidden=true}),
    g:addPolyline({{-5*vecK,5*vecK},{-5*vecI,4*vecI}}, {arrows=1, hidden=true})
)
g:Dballdots3d(Origin,nil,1.5)
g:Ddots3d({B1,C1, E1, F1},"gray")
g:Dpolyline3d({Origin, A, Cam},true,"dashed")
g:Dangle3d(Cam,A,Origin)
g:Dangle3d(Cam,Origin,v,0.3,"line width=0.8pt")
g:Darc3d(vecI,Origin,A,2.25,1,'-stealth'); g:Darc3d(vecK,Origin,Cam,2.25,1,'-stealth')
g:Dpolyline3d({{Cam,(B+Cam)/2},{Cam,C},{Cam,E},{Cam,F}},true,"->-=0.35,line width=0.1pt,gray")
g:Dpolyline3d({{B1,C1,E1,F1},{C1,F1}},true,"gray")
g:Dpolyline3d({B1,E1},true,"dotted,gray")
g:Ddots3d({B, C, E, F})
g:Dlabel3d("screen plane",D,{pos="NE",dir={u,v}})
g:Dlabel3d("target",Origin, {pos="S",dist=0.1}, "camera",Cam,{pos="SE"}, "$A$",B,{pos="N"},"$B$",C,{pos="S"},"$C$",E,{pos="N"},
"$D$",F,{},"$A'$",B1,{node_options="gray",dist=0}, "$B'$",C1,{pos="NW"}, "$C'$",E1,{pos="N"},"$D'$",F1,{}, "$\\theta$", 2.75*vecI,{pos="N",node_options="black"},"$\\varphi$", 1.3*(vecK+N), {}, "$z$",5*vecK,{},"$x$",4.5*vecI,{pos="center"})
local O = M(-3,-2,-4)
g:Ddots3d(O); g:Dpolyline3d({{O,O+vecI},{O,O+vecJ},{O,O+vecK}},'->')
g:Dlabel3d("$x$",O+1.25*vecI,{},"$y$",O+1.25*vecJ,{},"$z$",O+1.25*vecK,{},"Origin",O,{pos="S"})
u = g:Proj3dV(Cam)
g:Dpolyline3d( {Cam/2-0.4*vecK-Cam/3.5, Cam/2-0.4*vecK+Cam/3.5}, 'stealth-stealth')
g:Dlabel("$d$ = distance target - camera", g:Proj3d(Cam/2-0.4*vecK),{dir={u, cpx.I*u}, node_options="fill=white"})
dcamera(u,u)
g:Show()
\end{luadraw}
\captionof{figure}{Central projection}
\end{center}

The default values ​​are: \argu{theta}=$30$ (degrees), \argu{phi}=$60$, \argu{d}=$15$, \argu{target}=\varglob{pt3d.Origin} (the variable \varglob{pt3d.Origin} is the 3D point $(0,0,0)$).
